\chapter{Bond Futures}\label{ch:m2:bond-futures}

The Chicago Board of Trade's ten-year note future promises the delivery of a
note that does not exist. Its terms
describe a basket: any Treasury note with a remaining life between six and a
half and eight years at the start of the delivery month may be delivered, and
the seller, the short, chooses which one. To make notes with different
coupons and maturities interchangeable, each is scaled by a factor that
the exchange publishes; the scaling is exact only when yields are 6\%, which
they rarely are, so one note in the basket is always cheaper to deliver than
the others, and the future trades as a forward on that note plus the value of
the short's right to change its mind. Hedge funds hold hundreds of billions of
dollars of notes against short futures to earn the small difference between
the two. This chapter prices the contract, identifies the note it is really
written on, and values the choice.

\section{The contract and its basket}

\begin{definition}[Deliverable basket]\label{def:m2:bond-futures:basket}
The \emph{deliverable basket}\index{deliverable basket} of a bond future is
the set of securities the short may deliver, defined by the exchange's rules
by remaining (and sometimes original) maturity at the start of the delivery
month.
\end{definition}

\begin{dated}{2026-09}{Two ten-year contracts}\label{dat:m2:bond-futures:contracts}
\textbf{CBOT ten-year note future} (the ``6.5 to 8-year'' contract,
code \texttt{ZN}): USD~100\,000 face; tick half of a 32nd, USD~15.625; deliverable
notes with an original term of at most ten years and a remaining term of at
least six years six months and less than eight years; from the March 2026
contract month, reopened notes are deliverable too. Last trading day: the
seventh business day before the last business day of the delivery month;
delivery until the last business day. \textbf{Eurex Euro-Bund future}
(\texttt{FGBL}): EUR~100\,000 of a notional 6\% bond; deliverable German federal
bonds with 8.5 to 10.5 years remaining.
\end{dated}

\section{Conversion factors and invoice price}

\begin{definition}[Conversion factor and invoice price]\label{def:m2:bond-futures:cf}
The \emph{conversion factor}\index{conversion factor} of a deliverable note
is, approximately, the price of one unit of par at a yield of 6\%, computed by
the exchange's formula with the note's coupon rounded to the nearest eighth
and its remaining life from the first day of the delivery month rounded down,
to a whole quarter for the ten-year and bond contracts. A short that delivers
the note receives the \emph{invoice price}\index{invoice price}
\[
\text{invoice} \;=\; F \times \mathrm{CF} \;+\; \text{accrued interest at delivery},
\]
per 100 of face, where $F$ is the futures settlement price.
\end{definition}

\begin{method}[The exchange's conversion factor]\label{met:m2:bond-futures:cf}
With the coupon $c$ in decimals, $n$ whole years and $z$ whole months from the
first day of the delivery month to maturity ($z$ rounded down to a multiple of
3 for the ten-year and bond contracts), $v = z$ if $z < 7$ and $v = 3$
otherwise (for those contracts), set $a = 1.03^{-v/6}$, $b = \frac c2
\frac{6 - v}{6}$, $c' = 1.03^{-2n}$ if $z < 7$ and $1.03^{-(2n+1)}$ otherwise,
$d = \frac{c}{0.06}(1 - c')$, and
\[
\mathrm{CF} \;=\; a\,\Bigl[\frac c2 + c' + d\Bigr] - b,
\]
rounded to four decimals. It is the price at 6\% of a note whose next coupon
is $v$ months away, less the part of that coupon already accrued.
\end{method}

\begin{example}[The exchange's own check]\label{ex:m2:bond-futures:cf}
For the 3.75\% note of 15 November 2018 and the December 2008
contract, $n = 9$, the remaining 11 months are rounded down to $z = 9$, so $v =
3$: $a = 0.985329$, $b = 0.009375$, $c' = 0.570286$, $d = 0.268571$ and
$\mathrm{CF} = 0.8357$, the factor the exchange published. The chapter's build
reproduces the exchange's five worked examples to the fourth decimal.
\end{example}

The factor makes a note with a high coupon worth proportionally more
in delivery, so that every note in the basket would be equally attractive to
deliver if all yielded 6\%. At any other yield they are not. When yields are
below 6\% a long note gains more than the factor allows, so the short prefers
to deliver the note with the \emph{least} duration; when yields are above 6\%
it prefers the most.

\section{Cheapest-to-deliver}

\begin{definition}[Cheapest-to-deliver]\label{def:m2:bond-futures:ctd}
The \emph{cheapest-to-deliver}\index{cheapest-to-deliver} (CTD) is the note in
the basket that costs the short least to buy and deliver: the one with the
lowest price relative to what it will be invoiced at, equivalently, the
highest \omterm{def:m2:bond-futures:irr}{implied repo rate} (\cref{def:m2:bond-futures:irr}).
\end{definition}

\begin{proposition}[At delivery]\label{prop:m2:bond-futures:atdelivery}
On the last delivery day, if no arbitrage is possible, $F = \min_i
P_i/\mathrm{CF}_i$, where $P_i$ is the clean price of note $i$: the futures
price converges to the price of the cheapest note divided by its factor.
\end{proposition}
\begin{proof}
If $F > P_i/\mathrm{CF}_i$ for some $i$, sell the future, buy note $i$ and
deliver it: the invoice $F\,\mathrm{CF}_i + A_i$ exceeds the cost $P_i + A_i$. If
$F < \min_i P_i/\mathrm{CF}_i$, buy the future: the short must deliver some note
$j$ worth $P_j + A_j > F\,\mathrm{CF}_j + A_j$, which the long can sell for more
than it paid.
\end{proof}

Before delivery the future is priced off the CTD's forward, the price at which
buying the note today and financing it in repo breaks even at delivery, and
it trades a little below that, because the short owns the options of
\cref{sec:m2:bond-futures:options}.

\section{Gross basis, net basis and implied repo}

\begin{definition}[Gross basis, net basis, implied repo, basis trade]\label{def:m2:bond-futures:irr}
For a deliverable note with clean price $P$, factor $\mathrm{CF}$ and future
price $F$: the \emph{gross basis}\index{gross basis} is $P - F\,\mathrm{CF}$;
the \emph{net basis}\index{net basis} is the gross basis less the note's carry
to delivery (\cref{prop:m2:repo-and-specials:carry}); the \emph{implied repo
rate}\index{implied repo rate} is the financing rate at which buying the note,
financing it and delivering it into the future exactly breaks even,
\[
\text{IRR} \;=\; \frac{F\,\mathrm{CF} + A_{\text{del}} + C - (P + A_0)}{P + A_0}\,\frac{360}{d},
\]
with accrued interest $A_0$ today and $A_{\text{del}}$ at delivery, coupons
$C$ received in between and $d$ days. A \emph{basis trade}\index{basis trade}
buys the note and sells the future in the ratio of the factor (long the
basis), or the reverse (short the basis).
\end{definition}

The \omterm{def:m2:bond-futures:irr}{implied repo rate} is the futures market's quote for financing: a trader
who is long the basis earns it on the note's price and pays the actual \omterm{def:m2:repo-and-specials:rate}{repo
rate}, so the difference is the trade's profit if the note is delivered. For the
CTD it is usually a little below general collateral, and the gap, like the \omterm{def:m2:bond-futures:irr}{net
basis}, prices the short's options.

\begin{example}[A December basket]\label{ex:m2:bond-futures:basket}
On 25 September 2026 a December ten-year future trades at 111-24+
(111.765625); delivery is assumed on 31 December, 97 days away, and repo is
3.90\%. The table below shows five deliverable notes. The
3.875\% of August 2033, the shortest, has the smallest \omterm{def:m2:bond-futures:irr}{net basis},
two thirds of a 32nd, and the highest implied repo, 3.82\%, eight basis points
below repo: it is the CTD. The longest note would deliver at an implied repo
of 0.05\%. These notes and prices are illustrative.
\end{example}

\begin{center}
\small
\begin{tabular}{lrrrrr}
\toprule
Note & Clean price & Factor & \omterm{def:m2:bond-futures:irr}{Gross basis} & \omterm{def:m2:bond-futures:irr}{Net basis} & Implied repo \\
\midrule
3.875\% Aug-33 & 99.132 & 0.8870 & $-0.004$ & 0.020 & 3.82\% \\
4.125\% Nov-33 & 100.549 & 0.8971 & 0.284 & 0.261 & 2.95\% \\
4.25\% Feb-34 & 101.262 & 0.9012 & 0.539 & 0.488 & 2.12\% \\
4.625\% May-34 & 103.641 & 0.9200 & 0.817 & 0.695 & 1.45\% \\
4\% Aug-34 & 99.462 & 0.8806 & 1.041 & 1.036 & 0.05\% \\
\bottomrule
\end{tabular}
\end{center}

\begin{omfigure}
\begin{tikzpicture}[font=\small,
  box/.style={draw,thick,rounded corners=2pt,align=center,minimum height=0.95cm,minimum width=2.6cm,inner sep=4pt},
  arr/.style={-{Stealth[length=2.5mm]},thick}]
\node[box,draw=omThm,fill=omThm!8] (fund) at (0,0) {basis trader\\(long the basis)};
\node[box,draw=omDef,fill=omDef!8] (cash) at (-4.6,1.6) {cash market};
\node[box,draw=omProp,fill=omProp!8] (repo) at (-4.6,-1.6) {repo lender};
\node[box,draw=gray,fill=gray!8] (fut) at (4.6,0) {futures exchange\\(short future)};
\draw[arr] (cash) -- node[above=11pt,sloped,font=\footnotesize] {buys the CTD} (fund);
\draw[arr] (repo) -- node[below=11pt,sloped,font=\footnotesize] {finances it} (fund);
\draw[arr] (fund.north east) to[bend left=25] node[above=5pt,font=\footnotesize] {delivers the note} (fut.north west);
\draw[arr] (fut.south west) to[bend left=25] node[below=5pt,font=\footnotesize] {invoice $F\,\mathrm{CF} + A$} (fund.south east);
\end{tikzpicture}

\omcaption{The cash-and-carry \omterm{def:m2:bond-futures:irr}{basis trade}. The trader buys the cheapest note,
finances it in repo, sells futures in the ratio of the factor, and delivers
the note at expiry against the \omterm{def:m2:bond-futures:cf}{invoice price}. It earns the \omterm{def:m2:bond-futures:irr}{implied repo rate}
and pays the \omterm{def:m2:repo-and-specials:rate}{repo rate}; with leverage from repo, a few basis points on a large
position is the business.}\label{fig:m2:bond-futures:trade}
\end{omfigure}

\begin{dated}{2026-09}{The size of the basis trade}\label{dat:m2:bond-futures:basis}
A Federal Reserve note of June 2026 estimated aggregate Treasury cash-futures
basis positions at about USD~830 billion in September 2025, roughly twice the
early-2020 peak and 3.5\% of privately held Treasuries by market value.
Published estimates range from USD~350 billion to 1.5 trillion; hedge funds'
cash Treasury holdings reached USD~2 trillion at the end of 2025.
\end{dated}

\section{The delivery options}\label{sec:m2:bond-futures:options}

\begin{definition}[Quality option and wildcard option]\label{def:m2:bond-futures:options}
The short's right to choose which note to deliver is the \emph{quality
option}\index{quality option}. The \emph{wildcard option}\index{wildcard
option} is its right, during the delivery period, to decide whether to deliver
after the futures settlement price has been fixed for the day, while the notes
still trade; after the last trading day the \omterm{def:m2:bond-futures:cf}{invoice price} is fixed and the
short may deliver on any of the remaining days. The right to choose the day
within the delivery month is the \emph{timing} option.
\end{definition}

\begin{omfigure}
\begin{tikzpicture}[font=\footnotesize]
\draw[thick,-{Stealth[length=2.5mm]}] (0,0) -- (13.4,0) node[right]{time};
\draw[omDef,line width=3pt] (0.6,0.35) -- (9.0,0.35);
\foreach \x in {2.0,3.5,5.0,6.5,8.0}{\draw[omProp,thick] (\x,0.5) -- (\x,0.8);}
\node[anchor=south,text=omProp] at (5.0,0.85) {daily wildcard windows};
\node[anchor=south,text=omDef] at (4.8,1.35) {futures trade; settlement price fixed each day};
\draw[omThm,line width=3pt] (9.0,0.35) -- (12.6,0.35);
\node[anchor=south,text=omThm,align=center] at (10.8,0.5) {invoice fixed:\\deliver any day};
\foreach \x/\t in {0.6/first day of the month,9.0/last trading day,12.6/last delivery day}{
  \draw[thick] (\x,0.12) -- (\x,-0.12); \node[below=3pt] at (\x,-0.12) {\t};}
\end{tikzpicture}

\omcaption{The delivery month of the ten-year note future. The short chooses
the note (\omterm{def:m2:bond-futures:options}{quality option}), the day (timing option), and, every day until
trading stops, whether to deliver after the day's settlement price is known
(\omterm{def:m2:bond-futures:options}{wildcard option}). Trading stops on the seventh business day before the last
business day of the month; the \omterm{def:m2:bond-futures:cf}{invoice price} is then fixed, and delivery may
take place on any remaining business day up to the last.}\label{fig:m2:bond-futures:timeline}
\end{omfigure}

All three belong to the short, so the future is cheaper than the CTD's
forward by their value, and a trader long the basis, long the note and short
the future, owns them. The \omterm{def:m2:bond-futures:options}{quality option} is worth most when the CTD could
switch: when two notes' prices per factor are close, or when rates are
volatile enough to bring the switch within reach
(\cref{fig:m2:bond-futures:switch}).

\begin{omfigure}
\begin{tikzpicture}
\begin{axis}[width=11cm,height=5cm,
  xlabel={parallel change in yields (bp)},ylabel={price/factor minus cheapest},
  tick label style={font=\small},label style={font=\small},
  xmin=-200,xmax=250,ymin=-0.05,ymax=1.6,grid=major,grid style={gray!25},
  legend columns=3,legend style={font=\footnotesize,at={(0.5,-0.3)},anchor=north,draw=none,fill=none,/tikz/every even column/.append style={column sep=6pt}},
  table/col sep=comma]
\addplot[omDef,very thick] table[x=shift,y=d0]{figdata/markets-2/06-bond-futures/switch.csv};
\addplot[omProp,thick,dashed] table[x=shift,y=d1]{figdata/markets-2/06-bond-futures/switch.csv};
\addplot[gray,thick,dotted] table[x=shift,y=d2]{figdata/markets-2/06-bond-futures/switch.csv};
\addplot[omProp!60,thick] table[x=shift,y=d3]{figdata/markets-2/06-bond-futures/switch.csv};
\addplot[omThm,very thick] table[x=shift,y=d4]{figdata/markets-2/06-bond-futures/switch.csv};
\draw[omThm,thick,dashed] (axis cs:147,-0.05) -- (axis cs:147,1.6);
\node[font=\footnotesize,anchor=south west,text=omThm] at (axis cs:150,1.25) {switch: +147 bp};
\legend{Aug-33 (CTD today),Nov-33,Feb-34,May-34,Aug-34}
\end{axis}
\end{tikzpicture}

\omcaption{At delivery, each note's clean price divided by its factor, minus
the lowest, as all yields move together. The note on zero is the CTD. From
today's yields the shortest note stays cheapest until yields have risen 147
basis points, to about 5.5\%; beyond that the longest note takes over. Notes and
yields are illustrative. Data: the chapter's tutorial.}\label{fig:m2:bond-futures:switch}
\end{omfigure}

\begin{omfigure}
\begin{tikzpicture}
\begin{axis}[width=11cm,height=4.6cm,ybar,bar width=20pt,
  ylabel={implied repo (\%)},
  tick label style={font=\small},label style={font=\small},
  xmin=-0.6,xmax=4.6,ymin=0,ymax=4.6,
  xtick=data,xticklabels from table={figdata/markets-2/06-bond-futures/irr.csv}{label},
  xticklabel style={font=\footnotesize},
  table/col sep=comma]
\addplot[fill=omDef!55,draw=omDef] table[x=i,y=irr]{figdata/markets-2/06-bond-futures/irr.csv};
\draw[omThm,thick,dashed] (axis cs:-0.6,3.90) -- (axis cs:4.6,3.90);
\node[font=\footnotesize,anchor=south east,text=omThm] at (axis cs:4.55,3.92) {repo 3.90\%};
\end{axis}
\end{tikzpicture}

\omcaption{\omterm{def:m2:bond-futures:irr}{Implied repo rates} of the five notes of
\cref{ex:m2:bond-futures:basket}. The CTD's is highest and just below the \omterm{def:m2:repo-and-specials:rate}{repo
rate}; the others would lose money if bought and delivered. Data: the
chapter's tutorial.}\label{fig:m2:bond-futures:irr}
\end{omfigure}

\section{Tutorial: finding the cheapest-to-deliver}\label{tut:m2:bond-futures:ctd}

\begin{tutorial}
\textbf{Goal.} Compute the factors of a basket, its gross and net bases and
\omterm{def:m2:bond-futures:irr}{implied repo rates}, find the CTD, and locate the yield at which it switches.
\textbf{End state:} the table of \cref{ex:m2:bond-futures:basket} and
\cref{fig:m2:bond-futures:switch,fig:m2:bond-futures:irr}.
\begin{tutsteps}
\item \textbf{The factor}, from the exchange's formula.
\omcode{code/firm/ctd/firm_ctd.py}{23}{36}{The conversion factor.}\label{lst:m2:bond-futures:cf}
Check the exchange's examples: 0.9229, 0.8747, 0.8653, 0.8357, 0.7943.
\item \textbf{Basis and carry.}
\omcode{code/firm/ctd/firm_ctd.py}{53}{64}{Gross basis, carry to delivery and net basis.}\label{lst:m2:bond-futures:basis}
\item \textbf{Implied repo and the CTD.}
\omcode{code/firm/ctd/firm_ctd.py}{67}{76}{Implied repo rate, and the note that maximises it.}\label{lst:m2:bond-futures:irr}
\item \textbf{The switch.} \texttt{futures\_demo.switch\_shift()} scans
parallel shifts and reports the CTD changing at $+147$ basis points;
\texttt{fig\_futures.py} writes the two charts.
\end{tutsteps}
\textbf{What to change next.} Steepen the curve by adding a basis point per
year of maturity to each note's yield and find the new switch point; then
compute the value of the switch option for several volatilities
(\cref{exo:m2:bond-futures:7}).
\end{tutorial}

\section{Build: the basis analyser}\label{bld:m2:bond-futures:ctd}

\begin{build}
\bfield{Purpose} The miniature firm hedges its bond positions with futures
and runs \omterm{def:m2:bond-futures:irr}{basis trades}; both need the factor, the \omterm{def:m2:bond-futures:cf}{invoice price}, the basis and
the implied repo of every deliverable, every day.
\bfield{Interface} \texttt{conversion\_factor(coupon, first\_day, maturity,
quarter\_rounding)}; \texttt{invoice\_price(F, cf, accrued)};
\texttt{Deliverable(name, clean, accrued\_now, accrued\_delivery, cf,
coupon\_income)}; \texttt{gross\_basis}; \texttt{carry\_to\_delivery};
\texttt{net\_basis}; \texttt{implied\_repo}; \texttt{cheapest\_to\_deliver}.
\bfield{Rules} Factors by the exchange's formula; prices per 100; actual/360
financing; accrued interest from the bond library
(\cref{bld:m2:government-bonds:bond}).
\bfield{Acceptance tests} \texttt{code/firm/ctd/tests/}: the five published
factors; whole-month counting; a 6\% note maturing on the first day of the
delivery month has a factor of 1; the CTD's \omterm{def:m2:bond-futures:irr}{net basis} is zero at its own
implied repo and every other note's is positive.
\bfield{Stretch} The Eurex factor formula (annual coupons); delivery-date
choice from the timing option; a quality-option value from a curve model
rather than parallel shifts.
\end{build}

\begin{omsources}
\item CBOT, submission 25-099 to the CFTC (contract grade of the ten-year note
future), February 2025; CME Group, \emph{Calculating U.S. Treasury Futures
Conversion Factors}.
\item M.\ Henrard, \emph{Bond Futures: Description and Pricing}, OpenGamma
Quantitative Research, 2011.
\item Federal Reserve, FEDS Notes, ``Decomposing Hedge Funds' U.S. Treasury
Exposures'', June 2026; Office of Financial Research, ``Hedge Funds' Cash
Treasury Holdings Reach \$2 Trillion'', August 2026.
\item G.\ Burghardt and T.\ Belton, \emph{The Treasury Bond Basis}, 3rd ed.,
McGraw-Hill, 2005.
\end{omsources}

\section{Exercises}

\begin{exercise}[$\star$]\label{exo:m2:bond-futures:1}
The future settles at 111-24+. The CTD's factor is 0.8870 and its accrued
interest at delivery 1.453125. Give the \omterm{def:m2:bond-futures:cf}{invoice price} per 100 and per
contract.
\end{exercise}

\begin{exercise}[$\star$]\label{exo:m2:bond-futures:2}
The CTD's clean price is 99.132. Give its \omterm{def:m2:bond-futures:irr}{gross basis} in price and in 32nds.
Why can a \omterm{def:m2:bond-futures:irr}{gross basis} be negative?
\end{exercise}

\begin{exercise}[$\star$]\label{exo:m2:bond-futures:3}
Compute the December 2026 factor of a 4\% note maturing 15 November 2033, and
of a 6\% note with the same maturity. Why is the second not exactly 1?
\end{exercise}

\begin{exercise}[$\star\star$]\label{exo:m2:bond-futures:4}
The CTD's implied repo is 3.82\% and repo is 3.90\%. You buy the basis in
USD~100 million face and hold it to delivery with no switch. What do you lose?
Why might you still do it?
\end{exercise}

\begin{exercise}[$\star\star$]\label{exo:m2:bond-futures:5}
A fund holds USD~100 million face of the CTD and wants to hedge it with the
future. How many contracts should it sell? Why does the answer involve the
factor?
\end{exercise}

\begin{exercise}[$\star\star$]\label{exo:m2:bond-futures:6}
Explain from duration why the shortest note is cheapest to deliver when
yields are below 6\%, and why the switch goes to the longest note when yields
rise.
\end{exercise}

\begin{exercise}[$\star\star\star$]\label{exo:m2:bond-futures:7}
\emph{Coding.} With \texttt{futures\_demo.switch\_option\_value}, value the
switch option for a normal volatility of 80 basis points a year over the 97
days, then three times that. Report both in 64ths and comment.
\end{exercise}

\begin{exercise}[$\star\star\star$]\label{exo:m2:bond-futures:8}
\emph{Find the flaw.} ``The ten-year future trades at 111 while the notes it
delivers trade near 100: the future is 11\% expensive, sell it and buy the
notes.'' Correct the reasoning.
\end{exercise}

\section{Problem: The Switch}

\begin{problem}[{Weekend problem --- when the cheapest-to-deliver changes}]\label{pb:m2:bond-futures:1}
Take the basket and prices of \cref{ex:m2:bond-futures:basket}: the future at
111-24+, repo 3.90\%, 97 days to delivery.

\medskip\noindent\textbf{Part I --- The basket.}
\begin{enumerate}
\item Which notes are deliverable into the December contract, by the rule of
\cref{dat:m2:bond-futures:contracts}? Check each of the five.
\item Give the factor of the 4\% note of August 2034 and of the CTD.
\item Give the CTD's \omterm{def:m2:bond-futures:irr}{gross basis}, carry and \omterm{def:m2:bond-futures:irr}{net basis}.
\item Give its \omterm{def:m2:bond-futures:irr}{implied repo rate}. How far below repo is it?
\item Why is its carry negative?
\end{enumerate}

\medskip\noindent\textbf{Part II --- The switch.}
\begin{enumerate}[resume]
\item At delivery, with yields unchanged, what is each note's price per
factor above the CTD's?
\item By how much must all yields rise for the CTD to change? Which note takes
over?
\item What yield does the CTD have at that point?
\item Why is the switch below 6\%?
\item After the switch, how does the future's duration change?
\end{enumerate}

\medskip\noindent\textbf{Part III --- The option.}
\begin{enumerate}[resume]
\item With a normal volatility of 80 basis points a year, what is the standard
deviation of yields over the 97 days?
\item How many standard deviations away is the switch?
\item Give the value of the switch option at that volatility, in 64ths.
\item And at three times the volatility?
\item The future trades one 64th below the CTD's forward in the example. What
else could that 64th be paying for?
\end{enumerate}

\medskip\noindent\textbf{Part IV --- Judgement.}
\begin{enumerate}[resume]
\item A trader is long the basis of the CTD. Why is she long volatility?
\item Why does the switch matter to a hedger who uses futures, not only to a
basis trader?
\item What would a reopening of the Aug-33 note, adding to its supply, do to
the basis?
\item State the \emph{named result}: the shift at which the CTD switches, and
the value of the switch option at both volatilities.
\item In one sentence: what does a bond future really deliver?
\end{enumerate}
\end{problem}

\section{Interview questions}

\begin{interviewq}[$\star$ \iqroles{trader, bank}]\label{iq:m2:bond-futures:1}
What is a \omterm{def:m2:bond-futures:cf}{conversion factor}, and why does a Treasury future need one?
\end{interviewq}

\begin{interviewq}[$\star$ \iqroles{trader, researcher}]\label{iq:m2:bond-futures:2}
How do you identify the \omterm{def:m2:bond-futures:ctd}{cheapest-to-deliver}?
\end{interviewq}

\begin{interviewq}[$\star\star$ \iqroles{trader, researcher}]\label{iq:m2:bond-futures:3}
What is the \omterm{def:m2:bond-futures:irr}{implied repo rate}, and what does it mean if the CTD's implied repo
is above the actual \omterm{def:m2:repo-and-specials:rate}{repo rate}?
\end{interviewq}

\begin{interviewq}[$\star\star$ \iqroles{trader, bank}]\label{iq:m2:bond-futures:4}
Explain the Treasury \omterm{def:m2:bond-futures:irr}{basis trade} and what can go wrong with it.
\end{interviewq}

\begin{interviewq}[$\star\star$ \iqroles{researcher, developer}]\label{iq:m2:bond-futures:5}
How would you compute the DV01 of a Treasury future?
\end{interviewq}

\begin{interviewq}[$\star\star\star$ \iqroles{researcher, trader}]\label{iq:m2:bond-futures:6}
Which options does the short in a Treasury future hold, and how would you
value the most important one?
\end{interviewq}
